\section{Augmented Reality Technology}

GRAPHIC AR layers

2 layers: main info (HUD, doesn't move) + AR (hazard highlighting, moves)

HUD technically not AR, but extension of it

XR immersion situational awareness

inadequate tech puts movement in the image, flickering, glitching etc

XR opens up new avenues of interaction
not constrained to two dimensions -> can implement more "natural" movement, gestures

effects of motion on legibility

\subsection{Limitations of Optical See-Through Displays}

All industrial scenarios involve a certain level of danger or hazards, making it imperative to not cut off senses if system breaks, especially because the device cannot be taken off -> optical see-through. 

Regardless of used technology, images are formed with light
- additive color model, black is basically transparent \cite{kimEffectsDarkMode2019}
- competes with ambient lighting for visibility

Are there legal requirements/standards?
lighter color = more light + opacity -> eye strain?

\subsection{Color Blending}

\cite{gabbardPerceptualColorMatchingMethod2022} advises to use colors with high source luminance, with red usually not perceived as red -> would prefer to use yellow as warning color

Since Confined Spaces have minimal lighting (see section 1), they are an optimal background for readability. The brightness and opacity could be reduced. dark env + torch might impact readability

Industrial equipment, especially if it might be easily damaged, 
technology for augmented reality gets more expensive with either bigger field of view or better color rendering.
Making sure the proposed design foundations work in monochrome as well enables choosing more economic options, aside from accessibility.

---

\cite{itohIndistinguishableAugmentedReality2021}